\section{External validation and robustness of the ranking (Phase 3)}
\label{sec:external}

The ranking shipped above was built entirely from reference datasets (TCGA, GTEx,
HPA, DepMap, Open Targets). This section reports three post-lock audits that ask
whether the ranking survives contact with data it never saw, whether its safety
term is stable to a defensible re-weighting, and whether the composite score is
stable to perturbation of its locked weights. All criteria below were
pre-specified in committed files before the corresponding analysis ran
(\texttt{results/external\_bulk\_prespec.json}, then the analysis scripts and
outputs); no threshold was adjusted after the fact.

\subsection{Independent bulk tumor-versus-normal validation (GEO)}
Seven antigens from the locked candidate set were tested on independent GEO bulk
series with paired tumor and normal samples, one pre-specified series per
indication and one pre-specified probe per antigen (Table~\ref{tab:geo}). The
locked primary endpoint --- at least 6 of 7 antigens enriched in tumor at
Benjamini--Hochberg $q < 0.05$ --- \textbf{failed as measured}: 4 of 7 enriched
(CA9 $q = 5.0\times10^{-23}$, MSLN $q = 2.8\times10^{-6}$, SLC39A6
$q = 7.3\times10^{-5}$, CLDN6 $q = 9.5\times10^{-6}$), and CA12, PSCA, CLDN18 did
not enrich. The pre-specified secondary endpoint, rank concordance between the
external tumor medians and the rankscore tumor-signal term, was
Spearman $\rho = 0.179$ ($p = 0.70$), not significant, reported with no threshold
per the pre-specification.

\begin{table}[h]\centering\footnotesize
\caption{Independent GEO bulk validation, one series per indication, one
pre-specified probe per antigen. $q$ is Benjamini--Hochberg across the seven
tests. Medians are log-scale expression as distributed by each
series.}\label{tab:geo}
\begin{tabular}{llllrrrrl}
\toprule
Antigen & Series & Platform & Probe & $n_T$ & $n_N$ & Median $T$ & Median $N$ & Enriched \\
\midrule
CA9 & GSE53757 & GPL570 & 205199\_at & 72 & 72 & 11.63 & 6.71 & yes ($q=5.0{\times}10^{-23}$) \\
MSLN & GSE62165 & GPL13667 & 11719544\_a\_at & 118 & 13 & 6.98 & 4.87 & yes ($q=2.8{\times}10^{-6}$) \\
SLC39A6 & GSE42568 & GPL570 & 202088\_at & 104 & 17 & 12.67 & 10.89 & yes ($q=7.3{\times}10^{-5}$) \\
CLDN6 & GSE26712 & GPL96 & 208474\_at & 185 & 10 & 6.50 & 4.54 & yes ($q=9.5{\times}10^{-6}$) \\
PSCA & GSE46602 & GPL570 & 205319\_at & 36 & 14 & 5.30 & --- & no ($q=0.52$) \\
CA12 & GSE53757 & GPL570 & 203963\_at & 72 & 72 & 13.87 & 14.61 & no ($q\approx1$) \\
CLDN18 & GSE13861 & GPL6884 & ILMN\_1696284 & 65 & 19 & 12.39 & --- & no ($q\approx1$) \\
\bottomrule
\end{tabular}
\end{table}

The steering read, not a reframing: the two clearest non-enrichments are
informative in the direction of safety. CA12 median expression is \emph{higher}
in normal kidney than in tumor (14.61 vs 13.87), and CLDN18 shows no
tumor-over-normal window in gastric tissue --- exactly the two antigens the
safety models already rank at the bottom on both versions of the safety score
(below). A ranking component and an independent dataset pointing the same way is
stronger evidence than either alone; the primary endpoint's failure stands as
measured, and its information is used, not hidden.

\subsection{Safety model v2: equal-weight additive re-score}
The v1 safety term was a hand-weighted composite. Version 2 replaces it with an
equal-weight mean of the four min-max-normalized safety terms (normal-tissue
safety, essentiality, Mendelian-disease freedom, GWAS-burden freedom), computed
additively with v1 untouched. The v2 ordering is CLDN6 (0.998) $>$ CA9 (0.767)
$>$ SLC39A6 (0.748) $>$ MSLN (0.595) $>$ PSCA (0.542) $\approx$ CLDN18 (0.542)
$>$ CA12 (0.334), versus the v1 ordering SLC39A6 $>$ PSCA $>$ CLDN6 $>$ CA9 $>$
MSLN $>$ CLDN18 $>$ CA12; the two versions correlate at Spearman 0.607. The
cross-analysis coherence is the substantive point: CA12 and CLDN18 are bottom-two
on \emph{both} safety versions and are the two clearest GEO non-enrichments;
CLDN6 is best-or-near-best on all four safety terms and enriched externally.

\subsection{Weight-perturbation robustness of the locked ranking}
The locked five-term rankscore weights (tumor signal 0.30, normal safety 0.25,
protein breadth 0.15, tractability 0.15, non-essentiality 0.15) were perturbed
with 1{,}000 Dirichlet draws (concentration $=$ locked weights $\times$ 200, seed
20260929). The perturbed rankings agree with the locked ranking at median
Spearman 1.0 (IQR [0.964, 1.0], minimum 0.857); CA9 stays first in 100\% of
draws, the top-3 set holds in 99.7\%, and the top-5 set in 83.5\%. One-at-a-time
zero-out sensitivity: only zeroing normal safety or tractability changes the
podium; zeroing tumor signal, protein breadth, or non-essentiality leaves the
top three intact. The headline ranking is therefore a property of the evidence,
not of the exact weight vector.

\subsection{Field convergence: the CA9 AND-gate axis in 2026}
A 2026 preclinical construct, LNK001, pairs CAIX (CA9) with ENPP3 in a stringent
AND-gate dual-antigen CAR for clear-cell renal cell carcinoma --- independent
field validation of the axis this ranking puts first. Our committed bulk AND-gate
metric (\texttt{results/andgate\_p75\_igfiltered.json}, KIRC) ranks
CA9$\times$ENPP3 twelfth of fifteen pairs (gated window 2.886, gain 1.638); we
report that as measured. The mechanistic case for the gate is visible in our own
committed data: CA9's normal-tissue liability is concentrated in stomach (481
TPM, GTEx) and cerebellum (41--43 TPM), and ENPP3 is at or below 0.1 TPM in
stomach and 2.8 TPM in cerebellum, so a CA9 AND ENPP3 gate closes both windows
in bulk. From the same committed output we nominate CA9$\times$SLC17A3 as the
bulk-stronger KIRC pair (window 5.562, gain 6.179), with the caveat that
SLC17A3's normal kidney-cortex expression (33.6 TPM) limits its protection to
non-kidney tissue. The discriminating measurement --- single-cell co-occupancy
of the two antigens on the same malignant cells --- is queued as Tier-4 work;
bulk windows cannot settle it.
